\subsection{Replication of an isolated-cycle seed}\label{sec:seed-replication}
A generalized Li\'enard field \(X=(y-F(x),G(x))\) of degree at most \(n\) can be replicated at a calculable term cost. We first choose a common parameter giving enough hyperbolic cycles and then determine the support of the pullback.
\begin{proposition}[Replication of a finite seed]\label{prop:seed-strips}
Suppose \(X\) as above has \(K\ge1\) distinct selected isolated limit cycles.
For every integer \(k\ge1\) and every even integer \(t\ge2\), there exists a rational polynomial field with coprime components and pure univariate terms having at least \(tkK\) simultaneous hyperbolic limit cycles. Its term count and degree are at most
\[
 k+t(n+1),\qquad \max\{2k-1,t(n+1)-1\},
\]
respectively.
\end{proposition}

\begin{proof}
Consider \(X_a=(y-F(x)-aG(x),G(x))\). On a selected orbit of period \(T\) at \(a=0\), take an analytic transversal with coordinate \(s=0\) at the orbit and let \(d(s,a)\) be the return displacement. Its parameter derivative at \((0,0)\) is nonzero. Indeed, with the starting point fixed, the parameter variation \(w\) of the flow satisfies
\[
 \frac{\mathrm d}{\mathrm dt}\det(X,w)
   =(\operatorname{div}X)\det(X,w)+\det(X,\partial_aX_a|_0),
 \qquad \det(X,\partial_aX_a|_0)=G(x)^2 .
\]
The initial determinant is zero. The last term is positive somewhere on the orbit, since \(G(x)=\dot y\) cannot vanish identically on a nonconstant periodic orbit. Integration with the positive integrating factor gives a nonzero final determinant. The derivative of the hitting time adds a vector parallel to \(X\), and conversion to the transverse coordinate divides by a nonzero determinant.

The analytic zero set of \(d\) is therefore a graph \(a=A(s)\). Its order at zero is the finite multiplicity of \(d(s,0)\), since \(d(s,a)=(a-A(s))u(s,a)\) with \(u(0,0)\ne0\). After shrinking the section, \(A'(s)\ne0\) off zero. An odd-order graph meets each sufficiently small signed level once. An even-order graph meets levels of one sign twice and levels of the other sign zero times. At each intersection \(d_s=-uA'\ne0\), so the corresponding cycle is hyperbolic. Choose disjoint collars for the \(K\) cycles and one common bound on \(|a|\). If \(o\) cycles have odd order and \(e_+,e_-\) even-order cycles open to the positive and negative sides, the larger of the two local counts is
\[
 o+2\max(e_+,e_-)=K+|e_+-e_-|\ge K .
\]
Fix the favorable sign and choose \(K\) simple cycles of that one field. Rational approximation of \(F+aG\) and \(G\) within degree \(n\) preserves all of them. Write \(F,G\) for these rational polynomials.

Choose \(k\) distinct rational nodes \(b_j\) symmetric about zero. Hermite interpolation gives a unique polynomial \(\eta\) of degree at most \(2k-1\) with \(\eta(b_j)=0,\ \eta'(b_j)=1\) at every node. The data and uniqueness give rational coefficients and \(\eta(v)=-\eta(-v)\). Thus \(\eta\) has at most \(k\) monomials. Put \(H_R(v)=R\eta(v/R)\). In \(v=Rb_j+s\) one has \(H_R(Rb_j+s)=s+O(R^{-1})\) in \(C^1\) on each fixed compact \(s\)-interval. Consequently \((H_R(v)-F(x),G(x))\) converges in each translated chart to the hyperbolic seed. Take one sufficiently large integer \(R\) for the \(kK\) compact collars, with disjoint strips for different nodes.

Next take an integer \(C>0\) with the \(x\)-projections of these collars in \((-C,C)\). Put \(\phi(u)=CT_t(u)\), where \(T_t(\cos\theta)=\cos(t\theta)\). Each of the \(t\) monotone branches in \((-1,1)\) has a smooth inverse over the selected projections. The field
\begin{equation}\label{eq:seed-fold-field}
 Y=\bigl(H_R(v)-F(\phi(u)),\,\phi'(u)G(\phi(u))\bigr)
\end{equation}
satisfies \(D(\phi,\mathrm{id})Y=\phi'(u)(H_R(v)-F(\phi(u)),G(\phi(u)))\). On each branch its forward return is conjugate to the old return or to its inverse according to the sign of \(\phi'\). A positive multiplier different from one and its reciprocal are both hyperbolic. The \(tkK\) lifted cycles lie in distinct collars of this single field.

Since \(t\) is even, \(F(\phi)\) uses at most \(tn/2+1\) even powers of \(u\), including the constant, and \(\phi'G(\phi)\) uses at most \(t(n+1)/2\) odd powers. Adding \(k\) gives at most \(k+tn+t/2+1\le k+t(n+1)\), including \(t=2\). The two component degrees give the stated maximum. The second component is a nonzero polynomial in \(u\), and the highest \(v\)-coefficient in the first is a nonzero constant, so the pair is coprime.

For a cap \(\Delta\) and \(k+t(n+1)\le m\le(\Delta+1)(\Delta+2)\) with both degrees within the cap, Remark~\othernumber{rem:exact-fill}{1.4} of the main article gives exact \(m\) and rational coprime components. For every \(m\le3\Delta+2\) one may fill within pure-power and constant positions. These include \(v,\ldots,v^\Delta\) in the first component and \(1,u,\ldots,u^\Delta\) in both, giving \(3\Delta+2\) positions that include the original support. Adding small rational coefficients there preserves the displayed coprimality argument.

\end{proof}

For a published finite seed, Huzak and De Maesschalck give \(2\lfloor(n-2)/2\rfloor+\lfloor m/2\rfloor\) hyperbolic cycles for a field \((y-F(x),G(x))\) with \(\deg F=n\ge6,\ \deg G=m\ge2\) \cite[Theorem~1.1, p.~2]{HuzakDeMaesschalck2014}. Taking both degrees equal to \(12\) gives \(K=16\). With \(k=22,\ t=2\), Proposition~\ref{prop:seed-strips} gives \(2\cdot22\cdot16=704\) cycles using at most \(22+2(12+1)=48\) terms and degree at most \(\max(43,25)=43\). Filling vacant pure powers through degree \(43\) gives exactly \(48\) terms while preserving the cycles and coprimality. Thus this count applies under the diagonal cap \(47\).

\paragraph{A conditional comparison with linear restoring force.}
Assume the seed theorem of the unrefereed preprint \cite[Theorem~1.1]{LuClassical}. For each \(n\ge14\) it asserts that, with \(G(x)=-x\), an open set of degree-\(n\) primitives \(F\) gives at least \(K=\lceil(7n-48)/2\rceil\) simultaneous hyperbolic cycles. Choose a rational seed in that open set and an integer \(C\) whose open interval \((-C,C)\) contains the horizontal projections of the selected collars. A vertical translation replaces \(F\) by \(F-F(-C)\). Choosing \(R\) large in the Hermite construction keeps the translated copies inside the same horizontal range. For \(\phi(u)=C(2u^2-1)\), the first component of \eqref{eq:seed-fold-field} then occupies at most \(k\) odd \(v\)-powers and \(n\) positive even \(u\)-powers. The second is \(-\phi'\phi\), with just \(u,u^3\). Small rational changes fill these \(k+n+2\) positions, preserve the \(2kK\) cycles, and give degree \(\max(2k-1,2n,3)\). The leading \(v\)-coefficient is a nonzero constant, ensuring coprimality. With this seed the count under the diagonal cap is
\[
 n=\lfloor(N-1)/2\rfloor,\qquad k=N-n-2,\qquad
 \mathcal C(N)=2k\left\lceil(7n-48)/2\right\rceil\quad(n\ge14).
\]
On the grid \(N=2^j,\ 6\le j\le20\), the folded count \(\mathcal F(N)\) of the article, equation~(\paperlink{main}{equation.4.7}{eq:finite-main}{4.7}), first exceeds this quadratic count at \(2^{16}\). There \(\mathcal F(N)=7{,}547{,}387{,}904\) and \(\mathcal C(N)=7{,}514{,}193{,}974\).
